% SEGMENT OLP-0543-S01
% Part: set-theory
% Chapter: z
% Section: infinity-again

\documentclass[../../../include/open-logic-section]{subfiles}

\begin{document}

\olfileid{sth}{z}{infinity-again}
\olsection{முடிவிலி}

ஏதேனும் கணங்கள் இருந்தால், முடிவிலாப் பல கணங்கள் உள்ளன என்பதை உறுதிசெய்ய
நம்மிடம் ஏற்கெனவே போதுமான அடிகோள்கள் உள்ளன. ஏதோ ஒரு கணம் உள்ளது என்க;
அப்போது $\emptyset$ உள்ளது
(\olref[sth][z][sep]{prop:emptyexists} இன்படி). இப்போது எந்தக் கணம்~$x$
க்கும், \olref[sth][z][pairs]{prop:pairsconsequences} இன்படி
$x \cup \{x\}$ என்ற கணம் உள்ளது. ஆகவே இதனைச் சில முறை பயன்படுத்தினால்,
பின்வரும் கணங்கள் கிடைக்கும்:
\begin{enumerate}
	\item[0.] $\emptyset$ %& $0$ & படிநிலை $0$\\
	\item[1.] $ \{\emptyset\}$ %& $1$ & படிநிலை $1$\\
	\item[2.] $\{\emptyset, \{\emptyset\}\}$ %& $2$ & படிநிலை $2$\\
	\item[3.] $\{\emptyset, \{\emptyset\}, \{\emptyset, \{\emptyset\}\}\}$ %&$3$ & படிநிலை $3$\\
	\item[4.] $\{\emptyset, \{\emptyset\}, \{\emptyset, \{\emptyset\}\}, \{\emptyset, \{\emptyset\}, \{\emptyset, \{\emptyset\}\}\}\}$% & $4$ & படிநிலை $4$
\end{enumerate}
இந்தக் கணங்கள் ஒவ்வொன்றும் மற்றவற்றிலிருந்து வேறுபட்டது என்பதைச்
சரிபார்க்கலாம்.

சில காரணங்களுக்காக எண்மிடலை $0$ இலிருந்து தொடங்கியுள்ளோம். அவற்றுள் ஒரு
காரணம் இதுதான். ``$n$'' என்று நாம் பெயரிட்ட கணத்தில் துல்லியமாக $n$
உறுப்புகள் உள்ளன என்றும், அது (உள்ளுணர்வின்படி) $n$-ஆம் படிநிலையில்
உருவாகிறது என்றும் சரிபார்ப்பது கடினமல்ல.

ஆனால் இங்கு ஒரு சிக்கல் உண்டு. இது நமக்கு \emph{முடிவிலாப் பல} கணங்களைத்
தருகிறது; ஆனால் \emph{ஒரு முடிவிலிக் கணம்}, அதாவது முடிவிலாப் பல
உறுப்புகளைக் கொண்ட ஒரு கணம், உள்ளது என்பதை உறுதிசெய்யாது. இது உண்மையிலேயே
முக்கியமானது: ஒரு (டெடெகிண்ட்) முடிவிலிக் கணத்தைக் கண்டுபிடிக்க முடியாவிட்டால்,
ஒரு டெடெகிண்ட் இயற்கணிதத்தை நம்மால் கட்டமைக்க முடியாது. ஆனால் அதனை இயல்
எண்களின் கணமாகக் கருதுவதற்காக நமக்கு ஒரு டெடெகிண்ட் இயற்கணிதம் வேண்டும்.
(\olref[sfr][infinite][dedekindsproof]{sec} உடன் ஒப்பிடுக.)

முக்கியமாக, இதுவரை நாம் வகுத்த அடிகோள்கள் எந்த முடிவிலிக் கணத்தின்
இருப்பையும் \emph{உறுதிசெய்வதில்லை}. ஆகவே ஒரு புதிய அடிகோளை வகுக்க
வேண்டும்:

\begin{axiom}[முடிவிலி]
$\emptyset \in I$ ஆகவும், $x \in I$ ஆகும்போதெல்லாம்
$x \cup \{x\} \in I$ ஆகவும் உள்ள ஒரு கணம் $I$ உள்ளது.
\begin{align*}
	\exists I( & (\exists o \in I)\forall x\ x \notin o \land {}\\
	& (\forall x \in I)(\exists s \in I)\forall z(z \in s \liff (z \in x \lor z = x)))
\end{align*}
\end{axiom}

முடிவிலி அடிகோள் நமக்குத் தரும் கணம் $I$ டெடெகிண்ட்-முடிவிலியானது
என்பதை எளிதில் காணலாம். அதன் சிறப்பிக்கப்பட்ட உறுப்பு $\emptyset$;
$I$ மீதான ஒன்றுக்கொன்றான சார்பு $s(x) = x\cup \{x\}$ ஆல் தரப்படுகிறது.
இப்போது, ஒரு டெடெகிண்ட்-முடிவிலிக் கணத்திலிருந்து ஒரு டெடெகிண்ட்
இயற்கணிதத்தை எவ்வாறு பிரித்தெடுப்பது என்பதை
\olref[sfr][infinite][dedekind]{thm:DedekindInfiniteAlgebra} காட்டியது;
இதனை நமது இயல் எண்களின் கணமாகக் கருதுவோம். இன்னும் துல்லியமாக:
\begin{defn}\ollabel{defnomega}
முடிவிலி அடிகோள் நமக்குத் தரும் ஏதேனும் ஒரு கணம் $I$ ஐ எடுத்துக்கொள்க.
$s(x) = x \cup \{x\}$ என்ற சார்பை $s$ என்க. மேலும்
$\omega = \closureofunder{s}{\emptyset}$ என்க. $\omega$ இன் உறுப்புகளை
\emph{இயல் எண்கள்} என அழைக்கிறோம்; $\emptyset$ க்கு $s$ ஐ $n$ முறை
பயன்படுத்துவதால் கிடைக்கும் விளைவே $n$ என்று கூறுகிறோம்.
\end{defn}

சில பத்திகளுக்கு முன்பு ``$n$'' எனப் பெயரிடப்பட்ட கணமே $n$ என்ற எண்ணாகக்
\emph{கருதப்படும்} என்பதை இப்போது திரும்பிப் பார்த்துச் சரிபார்க்கலாம்.

இந்த நிர்ணயத்தின் முக்கியத்துவத்தை \olref[sth][z][nat]{sec} இல்
விவாதிப்போம். இப்போதைக்கு, அது ஓர் உள்ளுணர்வான முடிவை நிறுவ உதவுகிறது:

\begin{prop}\ollabel{naturalnumbersarentinfinite}
எந்த இயல் எண்ணும் டெடெகிண்ட்-முடிவிலியானதல்ல.
\end{prop}

\begin{proof}
இந்நிறுவல் தொகுத்தறிதலைப் பயன்படுத்துகிறது; அதாவது,
\olref[sfr][infinite][induction]{thm:dedinfiniteinduction} ஐப்
பயன்படுத்துகிறது. $0 = \emptyset$ டெடெகிண்ட்-முடிவிலியானதல்ல என்பது
தெளிவு. தொகுத்தறிதல் படிக்காக எதிர்மறைநேர்மாற்றை நிறுவுவோம்:
(முரணாக) $s(n)$ டெடெகிண்ட்-முடிவிலியானது எனில், $n$ உம்
டெடெகிண்ட்-முடிவிலியானது.

ஆகவே $s(n)$ டெடெகிண்ட்-முடிவிலியானது என்க; அதாவது,
$\ran{f}\subsetneq \dom{f} = s(n) = n \cup \{n\}$ ஆக அமையும் ஏதோ ஓர்
!!{injection} $f$ உள்ளது. இரண்டு நேர்வுகளைக் கருத வேண்டும்.

\emph{நேர்வு 1: $n \notin \ran{f}$.} ஆகவே $\ran{f} \subseteq n$,
மேலும் $f(n) \in n$. $g = \funrestrictionto{f}{n}$ என்க; இப்போது
$\ran{g} = \ran{f} \setminus \{f(n)\} \subsetneq n = \dom{g}$.
எனவே $n$ டெடெகிண்ட்-முடிவிலியானது.

\emph{நேர்வு 2: $n \in \ran{f}$.} ஏதோ ஒரு
$m \in \dom{f} \setminus \ran{f}$ ஐத் தேர்ந்தெடுத்து,
$s(n) = n \cup \{n\}$ என்ற பொருட்களத்தைக் கொண்ட $h$ என்ற சார்பைப்
பின்வருமாறு வரையறுக்கவும்:
\[
h(x) =
	\begin{cases}
		f(x) & f(x) \neq n\text{ எனில்}\\
		m & f(x)=n\text{ எனில்}
	\end{cases}
\]
ஆகவே $h$, $f$ ஆகியவை எல்லா இடங்களிலும் ஒன்றுபடுகின்றன; ஆனால்
$h(f^{-1}(n)) = m \neq n = f(f^{-1}(n))$ என்ற இடத்தில் மட்டும்
வேறுபடுகின்றன. $f$ ஓர் !!a{injection} என்பதால் $n \notin \ran{h}$;
மேலும் $\ran{h} \subsetneq \dom{h} = s(n)$. இப்போது நேர்வு 1 இன்
வாதத்தைப் பயன்படுத்தி, $n$ டெடெகிண்ட்-முடிவிலியானது எனப் பெறுகிறோம்.
\end{proof}

முடிவிலி அடிகோளை எவ்வாறு \emph{நியாயப்படுத்தலாம்} என்ற கேள்வி இன்னும்
எஞ்சுகிறது. நாம் சொல்லிவரும் கதையில் மேலும் ஒரு கொள்கையைச் சேர்க்க வேண்டும்
என்பதே சுருக்கமான பதில். அக்கொள்கை பின்வருவது:
\begin{enumerate}
	\item[] \stagesinf. ஒரு முடிவிலிப் படிநிலை உள்ளது. அதாவது, (a) முதல்
	படிநிலையாக இல்லாததும், (b) தனக்கு முந்திய சில படிநிலைகளைக் கொண்டதும்,
	ஆனால் (c) உடனடி முன்னோடி இல்லாததுமான ஒரு படிநிலை உள்ளது.
\end{enumerate}
இந்தக் கொள்கையிலிருந்து முடிவிலி அடிகோள் நேரடியாகப் பின்வருகிறது. இயல்
எண் $n$ படிநிலை $n$ இல் உருவாகிறது என்பதை அறிவோம். ஆகவே முதல் முடிவிலிப்
படிநிலையில் $\omega$ என்ற கணம் உருவாகிறது. மேலும் முடிவிலி அடிகோளுக்கு
$\omega$ தானே ஒரு சான்றாக அமைகிறது.

ஆனால் \stagesinf{} ஐ எவ்வாறு நியாயப்படுத்துவது என்ற கேள்விக்கே இது
நம்மை மீண்டும் கொண்டுசெல்கிறது. \stagessucc{} ஐப் போலவே, இதுவும்
திரள்வுறும் படிநிலை வரிசையைப் பற்றிய நமது கதையின் வெளிப்படையான பகுதியாக
இருக்கவில்லை. அதற்கும் மேலாக, கணங்கள் படிநிலை படிநிலையாக உருவாகும் ஒரு
படிநிலைமுறை வரிசை என்ற கருத்திலேயே, அச்செயல்முறை ஒரு \emph{முடிவிலிப்}
படிநிலையை உள்ளடக்க வேண்டும் என்று நம்மை எண்ணச் செய்வது எதுவுமில்லை.
படிநிலைகள் இயல் எண்களைப் போல வரிசைப்படுத்தப்பட்டுள்ளன என்று எண்ணுவது
முற்றிலும் இசைவானதாகத் தெரிகிறது.

ஆனால் இது ஓர் வெளிப்படையான சிக்கலை உருவாக்குகிறது.
\olref[sfr][infinite][dedekindsproof]{sec} இல், (எண்ணங்களின்) ஒரு
டெடெகிண்ட்-முடிவிலிக் கணம் உள்ளது என்பதற்கான டெடெகிண்டின் ``நிறுவலைக்''
கருதினோம். அது மிகவும் மனநிறைவளிப்பதாக உங்களுக்குத் தோன்றியிருக்காது.
ஆனால் கணங்களின் படிநிலைமுறைக் கருத்தாக்கம் (அல்லது ``சிந்தனையின்
விதிகள்'') \stagesinf{} ஐ ``நம்மீது கட்டாயப்படுத்தவில்லை'' எனில்,
ஒரு டெடெகிண்ட்-முடிவிலிக் கணம் உள்ளது என்ற கூற்றுக்கான உள்ளார்ந்த
நியாயம் இன்னும் நம்மிடம் இல்லை.

இதுபற்றி இன்னும் கூறுவதற்கு மிக அதிகம் உள்ளது. ஆனால் ஓர் அடிகோளை
(அல்லது கொள்கையை) நியாயப்படுத்துவதற்கு என்ன \emph{தேவைப்படலாம்} என்பதைப்
பற்றி நீங்கள் சிந்திக்கத் தொடங்கக்கூடிய நிலையை இப்போது அடைந்திருப்பீர்கள்
என்று நம்புகிறோம். தொடர்ந்து வருவதில் \stagesinf{} ஐ வெறுமனே
ஏற்றுக்கொள்வோம்.

\end{document}
